\documentclass[12pt]{article}
\usepackage{it_hw}
\input{author/macro}

\includeversion{problem}
\excludeversion{solution}

\begin{document}

\begin{center}
{\bf Information Theory and Its Application in Deep Learning: Midterm \begin{solution}Solutions\end{solution}} \\ \medskip
\begin{problem}Friday, April 26, 2024 \end{problem}
\end{center}

\begin{problem}
\noindent 
\begin{itemize}
\item You can use previous parts of a problem even if you did not solve them.
\item As throughout the course, entropy ($H$) and Mutual Information ($I$) are in bits.
\item $\log$ is taken in base 2. 
\item `prefix code' refers to a variable length code satisfying the prefix condition.  
\end{itemize}
\LeftFoot{Midterm}
\end{problem}

\begin{solution}
\LeftFoot{Midterm Solution}
\end{solution}

\begin{problem}

    \vskip .3in
    {\bf ID : \underline{\hspace{6cm}}}\\
    \vskip .3in
    {\bf Name : \underline{\hspace{5.5cm}}}\\
    ~
    \vskip 2in
    \begin{table}[h]
        \Large
        \begin{tabular}{c|c}
           1 &\textcolor{white}{aaaaaaaa}/~50\\
           2 &\textcolor{white}{aaaaaaaa}/~30 \\
           3 &\textcolor{white}{aaaaaaaa}/~30 \\
            \hline
           Total &  \textcolor{white}{aaaaaaaa}/~110 \\
        \end{tabular}
    \end{table}
    \newpage
    {\bf extra sheet}
    \newpage
\end{problem}

\begin{enumerate}


\item  {\bf Mix of Questions} {\it (50 points)}\\


\begin{enumerate}[1)]

% Entropy and mutual information
\item Let $Z_1, Z_2, Z_3, \ldots$ be i.i.d. random variables that take values ``0'' and ``1'' with equal probability. Further, let
\begin{equation}
X_i = \sum_{j=1}^i Z_j, \ \text{for} \ 1 \leq i \leq n.
\end{equation}

Find $I(X_1; X_2, X_3, \ldots, X_n)$.


% AEP
\item Let $U_1, U_2, U_3, \ldots$ be i.i.d. taking values $A, B, C, D, E$ and $F$, with the following distribution:
\begin{table}[htbp]
\begin{center}
\begin{tabular}{l|cccccc}
Symbol & $A$ & $B$ & $C$ & $D$ & $E$ & $F$ \\
Probability  & 1/2 & 1/4 & 1/8 & 1/16 & 1/32 &  1/32
\end{tabular}
\end{center}
\end{table}


\begin{enumerate}[(a)]
\item Compute $H(U_1)$.  
\item What is the most probable sequence of a given length $n$? What is its probability?
\item Recall the definition of the $\epsilon$-typical set for a memoryless source $U$: 
\begin{equation}
A_\epsilon^{(n)} = \left \{ u^n: \left| -\frac{1}{n} \log p(u^n) - H(U) \right| \leq \epsilon \right \}.
\end{equation} 
Does the sequence you found in part (b) belong to $A_\epsilon^{(n)}$ for $\epsilon= 0.1?$ How about for $\epsilon= 1$? 
\end{enumerate}

% AEP
\item Let $(X_i, Y_i)$ be i.i.d. $\sim p(x,y)$. Find the limit in probability, as $n \rightarrow \infty$, of 
\begin{equation}
\frac{1}{n} \log \frac{p(X^n, Y^n)}{p(X^n)p(Y^n)}.
\end{equation}



% Lossless coding
%\item Someone selects one of these symbols with the following probabilities:
%\begin{table}[htbp]
%\begin{center}
%\begin{tabular}{ll}
%Symbol & Probability\\
%
%A  & 1/2  \\
%B  & 1/4  \\
%C  & 1/8  \\
%D  & 1/16 \\
%E  & 1/32 \\
%F  & 1/32 \\
%
%\end{tabular}
%\end{center}
%\end{table}
%
%Your goal is to ask as \textbf{few} ``YES/NO'' questions as possible to find out the selected symbol. Which sequence of questions would you ask? Explain.


% Lossless coding
\item Consider a source with five symbols $u_1, u_2, u_3, u_4, u_5$, with probabilities $p(u_1) \geq p(u_2) \geq p(u_3) \geq p(u_4) \geq p(u_5)$.

\begin{enumerate}[(a)]
\item Suppose $p(u_1) \geq p(u_2) = p(u_3) = p(u_4) = p(u_5)$. Find the minimum value of $q$ such that $p(u_1) \geq q$ implies $n_1 = 1$. Here $n_1$ denotes the length of the codeword associated with symbol $u_1$ generated by a Huffman code applied to the source.

\item Suppose $p(u_1) \geq p(u_2) \geq p(u_3) \geq p(u_4) > p(u_5) = 0$. Find the largest value of $r$ such that $p(u_1) \leq r$ implies $n_1 > 1$. Here $n_1$ denotes the length of the codeword associated with symbol $u_1$ generated by a Huffman code applied to the source.

\end{enumerate}

% Lossless coding
\item Consider a random variable $X$ which takes on four possible values with probabilities $(1/3, 1/3, 1/4, 1/12)$.
\begin{enumerate}[(a)]
\item Construct a Huffman code for this random variable.
\item Show that there exist two different sets of optimal lengths for the codewords, namely, show that codeword length assignments $(1, 2, 3, 3)$ and $(2, 2, 2, 2)$ are both optimal.
\item Are there optimal codes with codeword lengths for some symbols that exceed the Shannon code length $\lceil \log \frac{1}{p(x)} \rceil$? (Hint: Check the codeword lengths from the previous part.)
\end{enumerate}



\end{enumerate}

%%% Solution of  Problem 1
%%%%%%%%%%%%%%%%%%%%%
\vskip 0.2in
\begin{solution}
{\bf Solution:}\\

\begin{enumerate}[1)]
\item 

\begin{align*}
I(X_1; X_2, X_3, \ldots, X_n) =& H(X_2,X_3, \ldots, X_n) - H(X_2, X_3, \ldots, X_n | X_1)  \\
=& H(X_2) + \sum_{i=3}^n H(X_i|X_2, \ldots, X_{i-1}) - \sum_{i=2}^n H(X_i|X_1, \ldots, X_{i-1})\\
=& H(X_2) + \sum_{i=3}^n H(Z_i) - \sum_{i=2}^n H(Z_i)\\
=& H(X_2) - H(Z_2)\\
=& \frac{3}{2} - 1\\
=& \frac{1}{2},
\end{align*}
since $X_2$ takes value ``0'' with probability $1/4$, value ``1'' with probability $1/2$ and value ``2'' with probability $1/4$, which gives $H(X_2) = 3/2$.

\item 
\begin{enumerate}[(a)]
\item $H(U_1) = 1/2 \log 2 + 1/4 \log 4 + 1/8 \log 8 + 1/16 \log 16 + 2/32 \log 32 = 31/16$.  
\item The most probable sequence is the symbol A repeated $n$ times, that is, AAAAAAA$\ldots$. \\
Its probability is $(\frac{1}{2})^n$.
\item Note that $-\frac{1}{n} \log p(u^n) = -\frac{1}{n} \log (\frac{1}{2})^n = 1$. We also have from part (a) that $H(U) = \frac{31}{16}$. Therefore,

\begin{align*}
\left| -\frac{1}{n} \log p(u^n) - H(U) \right| =& \left| 1 - \frac{31}{16} \right| \\
=& \frac{15}{16}\\
>& \ 0.1 \ \text{and} \ < 1
\end{align*}

Thus the most typical sequence $A^n$ belongs to $A_\epsilon^{(n)}$ for $\epsilon = 1$, but not for $\epsilon = 0.1$.



\end{enumerate}


\item 


First, note that we have $p(X^n, Y^n) = \prod_{i=1}^np(X_i, Y_i)$, $p(X^n) = \prod_{i=1}^np(X_i)$ and $p(Y^n) = \prod_{i=1}^np(Y_i)$.

\begin{align*}
\frac{1}{n} \log \frac{p(X^n, Y^n)}{p(X^n)p(Y^n)} &= \frac{1}{n} \log \frac{\prod_{i=1}^np(X_i, Y_i)}{\prod_{i=1}^np(X_i)\prod_{i=1}^np(Y_i)} \\
&= \frac{1}{n} \sum_{i=1}^n \log \frac{p(X_i, Y_i)}{p(X_i)p(Y_i)} \\
& \xrightarrow{n \rightarrow \infty} \mathbb{E}[\log \frac{p(X, Y)}{p(X)p(Y)}]\\
&=I(X;Y),
\end{align*}
where we have used the Law of Large Numbers.



\item 

\begin{enumerate}[(a)]
\item Note that $p(u_2) = p(u_3) = p(u_4) = p(u_5) = (1 - p(u_1))/4$. Without loss of generality, the Huffman code will first combine symbols $u_4$ and $u_5$, creating a super-symbol with probability $(1 - p(u_1))/2$, which is bigger than $(1 - p(1))/4$. Thus the Huffman code will then combine symbols $u_2$ and $u_3$, creating another super-symbol with probability $(1- p(u_1))/2$. For symbol $u_1$ to be assigned a codeword with length $1$, we need the Huffman code to combine next the two super-symbols. Therefore, we must have $p(u_1) > (1 - p(u_1))/2$. Or, equivalently,
\begin{equation}
p(u_1) > \frac{1}{3}.
\end{equation}



\item In this case, the Huffman code will first combine symbols $u_3$ and $u_4$ into a super-symbol with probability $p(u_3) + p(u_4)$. For the length of the codeword assigned to symbol $u_1$ to be bigger than $1$, we must have $p(u_1) < p(u_3) + p(u_4)$, which implies $p(u_2) < p(u_3) + p(u_4)$, since $p(u_2) > p(u_1)$. For $p(u_3) + p(u_4)$ to be as small as possible, we must have $p(u_1) = p(u_2)$, which implies $p(u_3) + p(u_4) = 1 - 2p(u_1)$. Thus $p(u_1)$ must satisfy $p(u_1) < 1 - 2p(u_1)$, which implies
\begin{equation}
p(u_1) < \frac{1}{3}.
\end{equation}


\end{enumerate}








\item 
\begin{enumerate}[(a)]
\item Applying the Huffman algorithm gives us the following table:
\begin{table}[htbp]
\begin{center}
\begin{tabular}{llllll}
Code & Symbol & Probability &  &  & \\
0    & 1 & 1/3  & 1/3 & 2/3 & 1 \\
11   & 2 & 1/3  & 1/3 & 1/3 &  \\
101  & 3 & 1/4  & 1/3 &  &  \\ 
100  & 4 & 1/12 &  &  &  \\
\end{tabular}
\end{center}
\end{table}
which gives codeword lengths of $1,2,3,3$ for the different codewords.


\item Both set of lengths $1,2,3,3$ and $2,2,2,2$ satisfy the Kraft inequality, and they both achieve the same expected length ($2$ bits) for the above distribution. Therefore they are both optimal.

\item The symbol with probability $1/4$ has an Huffman code of length $3$, which is greater than $\lceil \log \frac{1}{p(x)} \rceil$. Thus the Huffman code for a particular symbol may be longer than the Shannon code for that symbol. But on the average, the Huffman code cannot be longer than the Shannon code.

\end{enumerate}
\end{enumerate}

\end{solution}

\begin{problem}
\newpage {\bf problem 1 cont.}
\end{problem}
\begin{problem}
\newpage
\end{problem}

%%% Problem: Non-prefix Code
%%%%%%%%%%%%%%%%%%%%%
\item{\bf Non-prefix Code} {\it (30 points)}\\

Suppose $p(x)$ is a PMF over $\mathcal{X} = \{1,2,\cdots,K\}$, with $p(1) > p(2)>\cdots > p(K)$. We want to encode a random variable $X\sim p$. We care about encoding only this one random variable, therefore we do not require Unique Decodability but merely that the code be one-to-one, i.e., a different codeword for each of the $K$ source symbols. Note that even the zero length codeword is valid, i.e., sending nothing (the empty string) can represent one of the source symbols.
\begin{enumerate}[(a)]
% (a)
\item {\it (5 points)} Construct a coding scheme $c(X)$ that has the minimum expected code length. Let $l(i)$ be the length of the codeword of symbol $i$. Show that $l(i)= \lfloor \log i \rfloor$, where $\lfloor a \rfloor$ is the greatest integer no bigger than $a$.

% (b)
\item {\it (10 points)} Prove that the coding scheme from Part (a) satisfies  
\[l(i) \leq -\log p(i)\]
and conclude that the minimum expected code length is less than or equal to the entropy, i.e.,
\[\fE[l(X)]\leq H(X).\]

~[Hint : Note that $p(i)$ is the $i$-th largest value, and therefore, $P(i) \leq \frac{1}{i}$]

% (c)
\item {\it (15 points)} Show that
\[\fE[l(X)] \geq H(X)-1-\log(1+\log K ).\]
That is, lossless codes, even if not Uniquely Decodable, cannot beat the entropy by much.\\

~[Hint : You may want to use the fact that $\sum_{i=1}^K \frac{1}{i}\leq 1+\log K$]


\end{enumerate}

%%% Solution: Non-prefix Code
%%%%%%%%%%%%%%%%%%%%%
\vskip 0.2 in
\begin{solution}{\bf Solution: Non-prefix Code}\\
\begin{enumerate}[(a)]
% (a) sol
\item Assign the codewords in the following order
\[\phi,0,1,00,01,10,11,000,001,\cdots\]
This gives us $l(i)=\lfloor \log i\rfloor$.

% (b) sol
\item Note that $p(i)$ is the $i$-th largest value, therefore $p(i) \leq \frac{1}{i}$. Thus,
\begin{align*}
l(i) =& \lfloor \log i \rfloor\\
\leq &\log i\\
\leq &\log \frac{1}{p(i)}\\
=& -\log p(i)
\end{align*}
This implies $\fE[l(i)] \leq \fE[-\log p(X)] = H(X)$.

% (c) sol
\item 
\begin{align*}
\fE[l(X)] =&\sum_{i=1}^K p(i) \lfloor \log i \rfloor\\
\geq & \sum_{i=1}^K p(i) (\log i -1)\\
= & \sum_{i=1}^K p(i) (-\log p(i) + \log ip(i) -1)\\
= & H(X)-1-\sum_{i=1}^K p(i) \log \frac{1}{ip(i)}\\
\geq & H(X)-1- \log \left(\sum_{i=1}^K p(i)\frac{1}{ip(i)}\right)\\
\geq & H(X)-1- \log \left(\sum_{i=1}^K \frac{1}{i}\right)\\
\geq & H(X)-1- \log (1+\log K)
\end{align*}
\end{enumerate}
\end{solution}

\begin{problem}
\newpage {\bf problem 2 cont.}
\end{problem}
\begin{problem}
\newpage
\end{problem}



%%% Problem
%%%%%%%%%%%%%%%%%%%%%
\item  {\bf The prime number theorem} {\it (30 points)}\\

Some time around 300 B.C., someone showed that there are infinitely many prime numbers -- we know this because a proof appears in Euclid's famous {\em Elements}. In this problem, we will not only show that there are infinitely many prime numbers, but we will also give a lower bound on the rate of their growth using information theory. 
%\begin{definition}
%$\pi(n)$ is defined as the number of primes no greater than $n$. 
%\end{definition} 


Let $\pi(n)$ denote the number of primes no greater than $n$. Note that every positive   integer $n$ has a \textbf{unique} prime factorization of the form

\begin{align*}
n = \prod_{i=1}^{\pi(n)} p_{i}^{X_i},
\end{align*}
where $p_1, p_2, p_3, \ldots$ are the primes, that is, $p_1 = 2$, $p_2 = 3$, $p_3 = 5$, etc., and $X_i = X_i(n)$ is the non-negative integer representing the multiplicity of $p_i$ in the prime factorization of $n$. 

Let $N$ be uniformly distributed on $\{1, 2, 3 \ldots n\}$.

\begin{enumerate}
\item {\it (8 points)} Show that  $X_i(N)$ is an integer-valued random variable satisfying 
\begin{equation}
0 \leq X_i(N) \leq \log n.
\end{equation}

~[Hint : Try finding a lower and an upper bound for $p_i^{X_i(N)}$]

\item {\it (22 points)} Show that

\begin{equation}
\log n = H(N) \leq \pi(n) \log( \log n + 1).
\end{equation}

Thus, not only is $\pi(n) \rightarrow \infty$ but in fact  $\pi(n) \geq \frac{\log n}{\log (\log n + 1)} $.

~[Hint : Do $X_1(N), X_2(N), \ldots, X_{\pi(n)}(N)$ determine $N$? What does that say about the respective entropies?].  
\end{enumerate}

% Justify completely, each of the following steps:
%
%\begin{align}
%\log n &= H(N) \\
%&= H(X_1, X_2 \ldots X_{\pi(n)}) \\
%&\leq H(X_1) + H(X_2) + \ldots H(X_{\pi(n)}) \\
%&\leq \pi(n) \log( \log n + 1)
%\end{align}



% Hint: Use that $2^{X_i} \leq p_i^{X_i} \leq N \leq n$. 


%%% Solution of  Problem
%%%%%%%%%%%%%%%%%%%%%
\vskip 0.2in
\begin{solution}
{\bf Solution:}\\

\begin{enumerate}
\item $0 \leq X_i(N)$ is trivial. Note also that $2^{X_i} \leq p_i^{X_i} \leq N \leq n$. 

Thus, combining both results, $0 \leq X_i(N) \leq \log n$, as we wanted to show.

\item

\begin{align}
\log n &= H(N) \\
&= H(X_1, X_2 \ldots X_{\pi(n)}) \\
&= \sum_{i=1}^{\pi(n)} H(X_i|X_1, \ldots, X_{i-1})\\
&\leq H(X_1) + H(X_2) + \ldots H(X_{\pi(n)}) \\
&\leq \pi(n) \log( \log n + 1),
\end{align}
where the first step follows because there is a one-to-one mapping between $N$ and $X_1, X_2 \ldots X_{\pi(n)}$. The second step is by the chain rule for entropy. The next step is because conditioning reduces entropy, and the last one is because the distribution that maximizes entropy is the uniform one, there are $\pi(n)$ entropy terms, and the $X_i$'s can take at most $\log n + 1$ different values.

\end{enumerate}

\end{solution}

\begin{problem}
\newpage {\bf problem 3 cont.}
\newpage{\bf extrasheet.}
\end{problem}


\end{enumerate}

\end{document}
